\section{Background and Motivation}
\label{sec:bg_motivation}

We write this thesis at a moment when demand from AI and data-centric workloads is exploding. Modern computing is straining under simultaneous pressures of scale and efficiency: AI models are growing faster than improvements in transistor density and memory bandwidth, and as memory resides farther from compute, \emph{data movement} dominates both latency and energy. Traditional copper interconnects are fundamentally limited by RC delay and Ohmic loss: each bit transition requires charging and discharging long, resistive, capacitive lines (energy/bit $\propto CV^2$), and repeaters and equalization further increase latency and power \cite{Miller_InterconnectsReview,SunSiliconPhotonics}.

\paragraph{Optical interconnects and modulators.}
Optical interconnects address these limits by encoding electrical signals onto light via on-chip optical modulators and transmitting through low-loss waveguides. Free from $I^2R$ dissipation in the channel, such links sustain high bandwidth over millimeter–centimeter distances and, with wavelength-division multiplexing (WDM), achieve superior bandwidth density \cite{Miller_InterconnectsReview,SunSiliconPhotonics}. At the device level, modulation efficiency ultimately arises from a material’s \emph{nonlinear optical response}—field- or carrier-dependent changes in absorption and refractive index that convert small electrical drives into large optical contrast. In the classical regime, large modulation depth typically requires many photons or high optical fields. Pushing toward the regime where \emph{single} or \emph{few} photons can influence one another’s behavior demands much stronger interactions and carefully engineered light–matter coupling.

\paragraph{Toward quantum nonlinear optics.}
The long-term objective in quantum photonics is to realize solid-state platforms that are CMOS-compatible, uniform, and scalable, while achieving strong coupling to nanophotonic modes. Such platforms would enable few-photon phase shifts and non-Gaussian resources for quantum information processing, reducing the overheads inherent to purely probabilistic linear-optical schemes.

\paragraph{Why two-dimensional (2D) TMD semiconductors.}
Two-dimensional transition-metal dichalcogenides (TMDs) offer a compelling materials basis for both classical and quantum photonics. These van der Waals semiconductors are indirect-gap in bulk but become direct-gap in the monolayer, where tightly bound excitons dominate the optical response. Encapsulated monolayers exhibit exciton binding energies on the order of hundreds of meV, yielding robust resonances with large oscillator strengths and strong light–matter interaction at, and even above, room temperature. The broken inversion symmetry and strong spin–orbit coupling in monolayers produce valley-dependent optical selection rules, enabling helicity-selective excitation and potential valley-encoded information in classical and quantum settings.

\paragraph{Integration and device implications.}
The atomically thin form factor of TMDs supports heterogeneous, low-temperature back-end-of-line (BEOL) integration with silicon photonics and CMOS electronics. This facilitates precise placement of the active region at cavity field antinodes, minimizes parasitics to the electrical drivers, and enables compact, energy-efficient electro-optic modulators that leverage exciton-mediated changes in absorption and refractive index. Embedding TMDs in microcavities or photonic-crystal resonators can push devices into the strong-coupling regime (resolvable vacuum Rabi splitting), thereby enhancing per-photon nonlinear phase shifts and opening a pathway from high-contrast \emph{classical} modulation to \emph{few-photon} quantum nonlinear optics on a scalable, chip-integrated platform.
